% Adapted from the author's thesis, 05_Realisierung_neu.tex,
% fig:dialogverarbeitung_workcontext. Reflowed for one paper column.
% The implementation-specific context name is replaced by dialogue state.
\begin{tikzpicture}[
  font=\fontsize{8.5}{10}\selectfont, >=Latex,
  box/.style={draw,rounded corners=2pt,align=center,inner sep=4pt,
    minimum height=0.8cm,text width=2.5cm,fill=white},
  arrow/.style={->,thick},
  every node/.style={outer sep=1pt}
]
  \node[box,fill=blue!10] (message) at (0,0) {Current user input};
  \node[box] (state) at (3.9,0) {Previous\\dialogue state};
  \node[box,text width=3.5cm,minimum height=1cm] (process) at (1.95,-1.6)
    {Dialogue processing\\according to the selected\\orchestration variant};
  \node[box] (response) at (0,-3.2) {System response};
  \node[box] (updated) at (3.9,-3.2) {Updated\\dialogue state};
  \draw[arrow] (message.south) -- (process.north west);
  \draw[arrow] (state.south) -- (process.north east);
  \draw[arrow] (process.south west) -- (response.north);
  \draw[arrow] (process.south east) -- (updated.north);
  \draw[arrow] (updated.east) -- (5.7,-3.2) --
    node[right,rotate=90,anchor=south] {Next dialogue turn} (5.7,0) -- (state.east);
\end{tikzpicture}
